\subsection{Примеры комплексных сценариев}
\begin{codefile}{lua/examples/05_leds/leds_basic.lua}
\end{codefile}
\textbf{Пояснение.} В \verb|leds_basic.lua| создаётся объект \verb|Ledbar| и таблица \verb|action| с единственным состоянием \verb|"START"|, где по индексам выставляются разные цвета и в конце подаётся команда \verb|Ev.ENGINES_DISARM|. Функция \verb|callback(event)| по событиям \verb|Ev.ALTITUDE_REACHED|, \verb|Ev.POINT_REACHED| и \verb|Ev.COPTER_LANDED| вызывает соответствующее действие из таблицы состояний. Основной принцип --- связать индикацию с автоматом состояний и событиями автопилота. Это уже не базовое упражнение, а пример более цельного сценария.

\begin{codefile}{lua/examples/05_leds/leds_sleep.lua}
\end{codefile}
\textbf{Пояснение.} В \verb|leds_sleep.lua| используется тот же подход с таблицей \verb|action|, но между установками цвета вставлены вызовы \verb|sleep()|. Основной принцип --- показать, как последовательные блокирующие задержки влияют на выполнение сценария, и подготовить мотивацию для перехода к неблокирующим таймерам.
